\scanpage{179}
% The opening passage is enclosed and crossed by long diagonals.
\deleted{On peut supposer $X$ irréductible de pt générique $y$. L'inclusion $y\unsure{-|\to}X$ définit une suite spectrale (pour la cohomologie étale)
\[
H^*(y,\mathrm{G}_m)\Longleftarrow E_2^{pq}=H^p(X,\mathbb{R}^qi_{\unsure{*|x}}(\mathrm{G}_m))
\]
Or \deleted{\unsure{comme les}} comme les $H^q(y,\mathrm{G}_m)$ et les fibres de $\mathrm{R}^q i_*(\mathrm{G}_m)$, qui s'explicitent par de la cohomologie galoisienne, sont \deleted{\unsure{des}} groupes de torsion pour $q>0$, on trouve que \deleted{\unsure{l'on a} \unsure{sur|dans} les \unsure{les} groupes}
\[
H^p(X,\mathrm{R}_X^*)\tag*{$(*)$}
\]
\deleted{, où $\mathrm{R}_X=\mathrm{R}i_*(\mathrm{G}_m)$}\inserted{est un groupe de torsion pour $p>0$,} où $\mathrm{R}_X^*=\deleted{\unsure{\mathrm{R}i_*}}i_*(\mathrm{G}_{m\,\unsure{x|y}})$ est le faisceau (étale) des fonctions rationnelles \inserted{non nulles} sur $X$. \deleted{(Ceci étant}}

D'autre part, \inserted{si $X$ est \inserted{\unsure{noethérien}} ``sans idéaux premiers immergés'',} on a la suite exacte
\[
0\unsure{-|\to}\mathrm{G}_{m\,X}\longrightarrow\mathrm{R}_X^*\longrightarrow\underline{\mathrm{Div}}_X\longrightarrow0\tag*{(5.\unsure{1|2}.)}
\]
où $\underline{\mathrm{Div}}_X$ est le faisceau \inserted{(pour la \unsure{top|top.}} étale) des diviseurs (de Cartier) sur $X$, ce qui \deleted{\unsure{montre que}} \inserted{implique} compte tenu de $(*)$ que
\[
H^{i}(X,\underline{\mathrm{Div}}_X)\xrightarrow{\partial}H^{i+1}(X,\deleted{\unsure{\mathrm{R}|\mathrm{D}}}\mathrm{G}_{m\,X}).
\]
% slanted
\margin{\uline{donner le} \uline{corollaire} 5.4.}
est bijectif mod torsion \deleted{\unsure{si}} si $i\geq1$. \deleted{\unsure{Ainsi} \unsure{on aura|ne sera pas} vraiment utilisé} L'hypothèse de factorialité d'autre part implique \deleted{l'\unsure{isomorphisme}} qu'on a un isomorphisme de faisceaux étales :

% archive sheet number: 178
